\chapter{Conclusion}
ScamGuard spots scam and spam SMS aimed at Indian users and tells them why. Trained on our audited ScamGuard-IN-2026 corpus and tested with a campaign-aware split, the deployed detector gets 97.1\% of real messages from unseen campaigns right with 3.2\% false alarms, and it shows the exact phrases behind every verdict.

The cross-domain study that our guide proposed adds a second lesson. A classifier of the same kind, trained on another country's messages, loses about a third of its F1 on Indian SMS, and a small, carefully labelled local sample is enough to win it back. Anyone deploying such a filter should test it on local data first and label a few hundred local messages before turning to a heavier model.

All of this runs in a live application: quick check, full analysis, Hindi support, safety tools, a cross-domain page and a Streamlit demo. Every number in this report can be regenerated from the project repository.

\vspace{6pt}\noindent{\large\textbf{Future Enhancements}}
\begin{itemize}
\item Finish collecting the 2026 test set and repeat both evaluations on it.
\item Add a second annotator to the label audit and report their agreement.
\item Widen the cross-domain study to more target datasets and seeds, and try adaptation with less than 10\% of the pool.
\item Undo disguised spellings before features are built, and test against attackers who adapt or paraphrase.
\item Cover more Indian languages and compare small transformer models under the same protocols.
\end{itemize}
